%% ======================================================================
%% frontmatter/abstract.tex
%%
%% Abstract (and keywords).
%%
%% Structured as five paragraphs: (1) generic problem and gap,
%% (2) Polkadot justification and research question,
%% (3) system description and evaluation method, (4) results,
%% (5) three principal findings.
%%
%% Written in third person per CDLT guidelines (abstract must not use
%% first person). No citations or footnotes per guidelines.
%% Target ≤300 words; actual length approx. 295 words.
%% ======================================================================

\begin{abstract}

The content creator economy concentrates value in centralized intermediaries. No existing system, centralized or decentralized, provides subscription access, pay-per-view consumption, and permanent ownership through a single rights primitive that is portable across heterogeneous blockchain networks.

Polkadot's architecture is well-suited to investigate this gap. Its relay chain provides shared validator security, removing the external trust assumptions inherent in bridge-based solutions; XCM natively carries rights-management operations between parachains through a standardized protocol; and FRAME's composable module system allows subscription, pay-per-view, and permanent-ownership logic to coexist within a single runtime unit. This dissertation asks: \textit{how can a shared-security, multi-chain framework built on Polkadot XCM and FRAME primitives deliver a unified rights record natively supporting all three monetization models across heterogeneous blockchain networks, and what levels of transaction finality, cost efficiency, and creator revenue retention can it achieve?}

The Cross-Chain Content Rights Management Service (CCRMS) is a Polkadot parachain prototype designed to address this question. It has been evaluated through unit testing, performance benchmarks, a security audit, a centralized baseline comparison, and a Monte Carlo creator-revenue simulation.

The prototype achieves 26.2 transactions per second at a 96.9\% inclusion rate in the happy-path window of the stress test (Section~\ref{subsec:stress-test}), approximately six-second inclusion latency, and a destination-chain inclusion delay of up to seven blocks (median 5, approximately 30--35~s). The economic simulation yields approximately 16.8\% higher average creator revenue than centralized platforms, though this advantage is highly audience-dependent: creators above approximately 10,000 followers benefit decisively, whereas those below approximately 3,600 are better served by centralized discovery platforms. This result characterizes the CCRMS design; the benchmarked prototype bounds the number of rights holders per content item, which must be lifted before the advantage can be realized for larger audiences.

Three findings advance the state of knowledge. First, a unified rights record supporting all three access models can be built from standard FRAME primitives without novel cryptographic mechanisms, closing literature gap G7. Second, the revenue advantage of decentralized rights management is not universal, crossing a follower threshold at approximately 3,600. Third, conservative architectural choices provided resilience against an ecosystem-level disruption that occurred during the implementation period.

\end{abstract}

\noindent\textbf{Keywords:} cross-chain content rights, Polkadot XCM, unified rights record, creator revenue retention, design science research
